\section*{Chapitre \ref{ch:iv3:smart-contracts-and-crypto-engineering} --- Smart contracts et ingénierie crypto}
\markboth{\protect\omTocMark{}{Réponses~: Smart contracts et ingénierie crypto}}{\protect\omTocMark{}{Réponses~: Smart contracts et ingénierie crypto}}

\begin{solution}{iq:iv3:smart-contracts-and-crypto-engineering:1}
Les deux pièces entrent en entier. En sortent 600 pour le bénéficiaire et 190 sous forme d'une nouvelle pièce renvoyée à Alice, et les 10 manquants constituent les frais. Elle détient alors une seule pièce de 190. Sur une chaîne à comptes, son solde baisse simplement de 610, et celui du bénéficiaire augmente de 600.
\enlargethispage{-3\baselineskip}%
\omcode{code/interviews-3/08-smart-contracts-and-crypto-engineering/python/iv_utxo.py}{1}{17}{Paiement à partir de pièces.\protect\hfil\protect\break}\label{lst:iv3:cr:utxo}

Une pièce ne peut pas être dépensée partiellement, ce qui explique l'existence de la monnaie rendue. Une transaction prouve le droit de dépenser chaque entrée, et sa validité ne dépend que de ces entrées~; des transactions portant sur des pièces distinctes ne peuvent donc pas entrer en conflit, et elles peuvent être vérifiées en parallèle. Une double dépense consiste en deux transactions utilisant la même pièce, et une seule peut être acceptée.

Une chaîne à comptes tient à jour un solde et un nonce pour chaque adresse. Cela convient bien aux contrats, qui ont besoin d'un état partagé et modifiable, et cela permet des paiements simples et des transactions de petite taille. Cela impose aussi que les transactions d'un même compte soient ordonnées, et c'est à cela que sert le nonce.

Pour une bourse, ce choix affecte son infrastructure technique~: sur une chaîne à pièces, elle doit sélectionner les pièces et gérer les adresses de monnaie rendue, et sur une chaîne à comptes, elle doit gérer les nonces.

\iqwrong{«~La pièce de 500 paie 500, et 100 de la pièce de 300 paie le reste.~» Les pièces sont consommées en entier, et le reste revient sous forme de monnaie rendue.}
\end{solution}

\begin{solution}{iq:iv3:smart-contracts-and-crypto-engineering:2}
Les frais valent $21\,000 \times 21$ gwei $= 441\,000$ gwei, soit 0.000441 ETH. La part correspondant aux frais de base, 0.00042 ETH, est brûlée, et les frais de priorité, 0.000021 ETH, reviennent au producteur du bloc. Après un bloc plein, les frais de base suivants peuvent augmenter de 1/8, soit 12.5~\%, ici jusqu'à 22.5 gwei.

Le gas mesure un travail, et ce n'est pas une monnaie. Le prix du gas varie avec la demande, si bien que le même transfert peut coûter beaucoup plus cher à une heure d'affluence, tout en consommant les mêmes 21\,000 gas. L'expéditeur indique un plafond de frais et un pourboire, et paie les frais de base au moment de l'inclusion ainsi que le pourboire, dans la limite de ce plafond, le solde restant lui revenant.

La cible d'un bloc est la moitié de sa limite. Au-delà de la cible, les frais de base augmentent, de 1/8 au plus pour un bloc entièrement plein, et en deçà ils baissent. Ainsi, une hausse soudaine de la demande est intégrée au prix sur plusieurs blocs, et le pourboire détermine l'ordre au sein d'un bloc lorsque beaucoup veulent y entrer.

Un système de trading traite donc le gas comme un coût réel et volatil. Une transaction qui échoue paie néanmoins le gas qu'elle a consommé.

\iqwrong{«~La totalité des frais va au validateur.~» Avec cette règle, les frais de base sont brûlés, et seul le pourboire parvient au producteur du bloc.}
\end{solution}

\begin{solution}{iq:iv3:smart-contracts-and-crypto-engineering:3}
Ces appels renvoient 0, 1, 3e18 et 0. \texttt{divFirst} tronque $1/3$ à 0 avant la multiplication. Pour $(2^{255}, 2, 2)$, \texttt{mulFirst} échoue (revert), car $a \times c$ déborde sur 256 bits avant d'obtenir un résultat qui tient. Solidity ne dispose pas de virgule flottante exploitable, car tous les nœuds doivent obtenir des résultats identiques, et seuls les entiers le garantissent. On écrit 1.5 comme $1.5 \times 10^{18}$.

Les montants sont donc des entiers dans la plus petite unité~: 1 ETH vaut $10^{18}$ wei. La forme courante est la virgule fixe à 18 décimales, où \texttt{wmul} multiplie puis divise par $10^{18}$ pour rétablir l'échelle. \texttt{wmul(1, 1)} met la perte en évidence~: le résultat exact, $10^{-18}$ d'unité, s'arrondit à 0.

La règle consiste à multiplier avant de diviser, pour conserver la précision, et à surveiller le dépassement de capacité au milieu du calcul. À partir de la version 0.8, un dépassement annule la transaction~: un grand résultat intermédiaire conduit donc à une transaction échouée et non à un nombre faux. Une multiplication suivie d'une division avec un intermédiaire complet sur 512 bits constitue la parade.

Toute division arrondit aussi, et il faut décider au profit de qui. Un contrat doit en général arrondir au détriment de l'utilisateur, sinon de petites erreurs d'arrondi répétées peuvent se transformer en une fuite lente.

\iqwrong{«~Utilisez un type comme \texttt{fixed128x18}.~» Ces types peuvent être déclarés, sans toutefois pouvoir encore recevoir ni fournir de valeur, si bien qu'on utilise en pratique des bibliothèques bâties sur les entiers.}
\end{solution}

\begin{solution}{iq:iv3:smart-contracts-and-crypto-engineering:4}
\texttt{viaMemory()} renvoie 1, puis \texttt{viaStorage()} renvoie 2. Une variable \texttt{memory} reçoit une copie lorsqu'on lui affecte une donnée de stockage, si bien que la modification est perdue à la fin de l'appel. Une variable \texttt{storage} est une référence, de sorte qu'y écrire modifie l'état du contrat. Le stockage est de loin le plus coûteux.

Les trois emplacements diffèrent par leur durée de vie. \texttt{storage} est l'état persistant du contrat sur la chaîne, qui subsiste jusqu'à modification. \texttt{memory} est un espace de travail temporaire, valable pour l'appel externe en cours. \texttt{calldata} est l'entrée en lecture seule de l'appel. Utiliser \texttt{calldata} pour un argument de type tableau évite une copie, si bien que c'est souvent le moins coûteux.

Le piège, ici, est que \texttt{Point memory m = p} ressemble à une poignée sur \texttt{p}. Une fonction qui met à jour une position par l'intermédiaire d'une telle copie ne fait rien, sans aucune erreur. Elle peut aussi coûter peu de gas, ce qu'on peut prendre pour une bonne nouvelle. Et le bogue inverse existe~: une copie supposée qui est en réalité une référence vers le stockage, de sorte qu'une modification censée être temporaire est enregistrée.

\iqwrong{«~Les deux renvoient 2.~» Cela traite \texttt{memory} comme s'il s'agissait d'un alias. Dans ce sens, le compilateur effectue une copie.}
\end{solution}

\begin{solution}{iq:iv3:smart-contracts-and-crypto-engineering:5}
Le nonce est un compteur par compte, et une transaction doit porter la valeur suivante. Cela empêche une rediffusion~: un transfert signé avec le nonce 7 devient invalide dès que 7 est utilisé. Les nonces 0 et 1 peuvent être inclus, et 3 attend l'apparition de 2. Parmi deux transactions portant le même nonce, une seule au plus peut être incluse.

Le nonce impose donc un ordre strict par compte, et chaque transaction signée ne peut produire son effet qu'une seule fois. Comme la signature couvre le nonce, personne ne peut le modifier sans invalider la signature.

L'attente a des conséquences. Un nonce manquant bloque tout ce qui le suit, même avec des frais élevés. Un expéditeur qui émet de nombreuses transactions doit donc suivre lui-même ses nonces et combler rapidement les trous.

Deux transactions portant le même nonce constituent la manière habituelle de remplacer une transaction en attente. On envoie le même nonce avec des frais plus élevés, et la plupart des nœuds acceptent ce remplacement. Celle qui est incluse la première l'emporte, et l'autre ne peut plus jamais être incluse. On ne peut donc annuler un paiement en attente qu'en faisant passer un remplacement en premier, par exemple un transfert de valeur nulle vers son propre compte, et rien n'assure que cela fonctionne.

\iqwrong{«~Un nonce rend une transaction privée ou aléatoire.~» Ici, ce n'est qu'un numéro de séquence, visible de tous.}
\end{solution}

\begin{solution}{iq:iv3:smart-contracts-and-crypto-engineering:6}
Un bloc vu à l'instant peut être remplacé par une autre chaîne, si bien qu'un dépôt crédité pourrait disparaître après que l'utilisateur a retiré ou échangé ses fonds. Deux époques de $12 \times 32 = 384$ secondes chacune font 12.8 minutes, et la documentation d'Ethereum cite environ 15. Remettre en cause un bloc finalisé exige qu'un attaquant renonce à au moins un tiers de l'ensemble des ETH mis en jeu.

Un bloc situé à un seul créneau de profondeur ne bénéficie d'aucune protection de ce type. Il peut être évincé lors d'une réorganisation courte ordinaire, sans attaque et sans coût pour quiconque. Le coût d'une remise en cause croît avec la profondeur, et avec la preuve d'enjeu (proof of stake) il bondit dès que le bloc est finalisé.

Une bourse choisit donc une profondeur par chaîne et par montant, en arbitrant entre le délai imposé à l'utilisateur et le risque encouru. Sur une chaîne à preuve de travail (proof of work), il n'existe pas de point final~: la sécurité croît avec la profondeur, et une chaîne à faible puissance de hachage est peu coûteuse à attaquer, de sorte qu'elle exige davantage de confirmations. Pour une chaîne dotée d'un séquenceur rapide sur une couche 2, une confirmation provisoire est la promesse d'un seul acteur, tandis que l'état définitif sur la couche 1 constitue la réalité, qui peut n'intervenir que bien plus tard.

En interne, on peut créditer le solde de trading avant le solde retirable. Cela permet à l'utilisateur de négocier, mais les fonds ne peuvent sortir qu'une fois définitifs, de sorte qu'une remise en cause coûte au plus une position récupérable.

\iqwrong{«~Elle est dans un bloc, donc elle a eu lieu.~» Un bloc proche du sommet de la chaîne est une proposition qui survit le plus souvent.}
\end{solution}

\begin{solution}{iq:iv3:smart-contracts-and-crypto-engineering:7}
Une première écriture, de zéro vers une valeur non nulle, coûte 22\,100 gas. Une écriture d'une valeur non nulle vers une autre coûte 5000. Réécrire la même valeur coûte 2200. Une écriture vers zéro coûte 5000, et elle donne droit à un remboursement de 4800. Nous avons mesuré ces valeurs sur l'EVM local, déduction faite du surcoût de 392 gas lié à l'appel.

Chaque coût a deux composantes. Le premier accès à un emplacement dans une transaction est dit froid, à 2100, et les suivants sont dits chauds. Vient ensuite l'écriture~: 20\,000 pour faire naître un nouvel emplacement, puisque chaque nœud doit alors le conserver définitivement~; 2900 pour modifier un emplacement existant~; 100 si rien ne change. Soit $22\,100 = 20\,000 + 2100$, $5000 = 2900 + 2100$ et $2200 = 100 + 2100$.

Le stockage est tarifé ainsi, car l'état est une charge permanente pour chaque nœud, alors que le calcul ne fait que transiter. Vider un emplacement donne donc droit à un remboursement, puisqu'il allège cette charge. Les remboursements sont plafonnés au cinquième du gas consommé dans une transaction.

Dans une même transaction, un emplacement déjà touché est chaud, et le réécrire coûte bien moins~: c'est pourquoi le programme envoie des transactions distinctes, car à l'intérieur d'une seule, les écritures ultérieures trouveraient l'emplacement chaud et les chiffres seraient différents. Il faut donc regrouper ses écritures et lire un emplacement dans une variable locale si on en a besoin plus d'une fois, et il ne faut pas stocker ce qu'on peut émettre sous forme d'événement~: les logs coûtent bien moins cher, bien que les contrats ne puissent pas les lire.

\iqwrong{«~Écrire un mot coûte un montant fixe.~» Le coût dépend de l'ancienne valeur, de la nouvelle, et de savoir si l'emplacement a déjà été touché.}
\end{solution}

\begin{solution}{iq:iv3:smart-contracts-and-crypto-engineering:8}
\texttt{Loose} utilise 3 emplacements, et \texttt{Tight} en utilise 2. Dans \texttt{Tight}, l'emplacement 0 contient \texttt{a} dans ses 16 octets de poids faible et \texttt{c} dans ses 16 de poids fort, soit, à la lecture, \texttt{0x03}, des zéros, puis \texttt{01}. L'économie au déploiement correspond à un emplacement neuf, soit environ 22\,000 gas d'après le coût d'une écriture depuis zéro. Nous avons lu les emplacements directement sur l'EVM local, sans mesurer les déploiements eux-mêmes.

La règle consiste à remplir les emplacements dans l'ordre de déclaration. Une variable qui ne tient pas dans la partie restante de l'emplacement courant en ouvre un nouveau. Un \texttt{uint256} occupe un emplacement entier, si bien que dans \texttt{Loose} il sépare les deux variables \texttt{uint128}.

Le compactage a un coût et un avantage au-delà du déploiement. Deux variables dans un même emplacement s'écrivent en une seule opération si on les met à jour ensemble, et lire l'une rend l'autre emplacement chaud. Mais si l'une change souvent et l'autre non, chaque mise à jour doit lire, masquer et réécrire l'emplacement, et les emplacements partagés peuvent aussi rendre le code moins lisible. Il faut donc regrouper les variables selon ce qui est utilisé ensemble, et non seulement selon leur taille.

Dans un contrat évolutif, la disposition du stockage fait partie des données. L'état déployé reste en place, et une nouvelle implémentation lit les mêmes emplacements par leur position. Réordonner fait donc lire à la nouvelle version \texttt{b} là où se trouvait \texttt{a}, sans qu'aucune erreur ne soit signalée. Il ne reste qu'à ajouter des variables en fin de liste.

\iqwrong{«~Le compilateur réordonne les variables automatiquement.~» Il conserve l'ordre de déclaration, car les contrats s'appuient sur cette disposition.}
\end{solution}

\begin{solution}{iq:iv3:smart-contracts-and-crypto-engineering:9}
\texttt{checkedDec(0)} échoue (revert) avec le code de panique 0x11. \texttt{uncheckedDec(0)} s'enroule jusqu'à $2^{256} - 1$. \texttt{narrow(300)} renvoie 44, car une conversion explicite tronque simplement les bits de poids fort, sans aucune vérification. Ainsi, l'arithmétique est protégée par défaut depuis la version 0.8, sauf si l'on y renonce explicitement, mais les conversions ne le sont pas.

Avant la version 0.8, toute l'arithmétique s'enroulait, si bien qu'une soustraction passant sous zéro pouvait donner un solde énorme. Les vérifications coûtent désormais un peu de gas à chaque opération.

\texttt{unchecked} existe pour éviter ce coût là où la vérification ne peut pas échouer. Un compteur de boucle borné par la longueur d'un tableau ne peut pas déborder. Une différence qu'on vient de comparer est sûre~: si l'on a vérifié $a \ge b$, alors $a - b$ ne pose pas de problème. Le code de hachage et celui qui calcule des différences de temps peuvent au contraire vouloir un enroulement, et ce délibérément.

Il n'est pas sûr dès qu'une valeur fournie par un utilisateur intervient. Dans un transfert de jetons sans contrôle du solde, \texttt{unchecked \{ balance -= amount; \}} reproduit intégralement l'ancien bogue.

La conversion est le point à surveiller lors de la relecture. Une valeur issue d'un utilisateur, convertie en \texttt{uint96} pour gagner de la place, peut passer toutes ses vérifications sur 256 bits, puis perdre ses bits de poids fort au moment de son stockage. Il faut recourir à une conversion vérifiée issue d'une bibliothèque, ou à une comparaison explicite, à chaque réduction de largeur.

\iqwrong{«~La version 0.8 a corrigé le débordement.~» Elle a corrigé les opérateurs arithmétiques hors de \texttt{unchecked}. Les conversions et l'assembleur n'ont jamais eu cette vérification.}
\end{solution}

\begin{solution}{iq:iv3:smart-contracts-and-crypto-engineering:10}
Le vault finit à 0, et l'attaquant à 11 ETH, ce qu'indique notre exécution locale. Le retrait de Bob échoue ensuite. Le défaut tient à ce que \texttt{withdraw} envoie les fonds avant de remettre le solde à zéro. Il faut corriger le défaut en mettant d'abord le solde à zéro (vérifications, effets, interactions), ou en ajoutant un verrou. Mettre le solde à zéro en premier, et ajouter aussi un verrou là où les appels sont complexes.

Étape par étape~: l'attaquant dépose 1 et appelle \texttt{withdraw}. Le vault lit un solde de 1 et l'envoie. L'envoi vers un contrat exécute son \texttt{receive}, qui rappelle \texttt{withdraw}. Le solde vaut toujours 1, car la ligne qui le remet à zéro n'a pas encore été atteinte, si bien que le vault paie de nouveau. Cela se répète tant que le vault conserve au moins 1 ETH. Ce n'est qu'au retour des appels que les écritures à zéro sont effectuées, toutes trop tard.

Le code présenté ne comporte d'erreur dans aucune de ses lignes prise isolément. Le défaut tient à ce qu'on cède le contrôle à un contrat inconnu alors que l'état n'est pas à jour. La documentation de Solidity indique qu'un appel externe doit être la toute dernière étape d'une fonction.

Nous avons exécuté les deux correctifs contre le même attaquant. La ligne de remise à zéro étant remontée, le \texttt{withdraw} interne ne trouve rien et échoue (revert), ce qui annule toute l'attaque, et Bob retire ses 10. Avec un verrou sur la fonction, l'appel interne échoue (revert) sur le verrou, avec le même résultat. Le premier correctif agit à la racine et ne coûte aucun stockage. Un verrou protège aussi contre des chemins auxquels on n'a pas pensé, entre fonctions, bien qu'un verrou sur une seule fonction ne voie pas une réentrance passant par une autre. Le vault ici n'est qu'un contrat local de test, et rien de plus.

\iqwrong{«~Utiliser \texttt{transfer}, puisque sa limite de gas empêche la réentrance.~» Cela suppose que les coûts en gas restent inchangés, et les mises à niveau les ont modifiés.}
\end{solution}

\begin{solution}{iq:iv3:smart-contracts-and-crypto-engineering:11}
\texttt{paid} affiche true, le bénéficiaire détient 0, et \texttt{Payout} garde 1 ETH~: l'appel a échoué, et personne ne l'a remarqué. \texttt{payChecked} échoue (revert), si bien que rien n'est enregistré. Les appels de bas niveau renvoient un indicateur et n'échouent pas, et la version sans contrôle l'ignore, ce que le compilateur signale. Payer dans une boucle permet à un seul bénéficiaire de bloquer tous les autres.

Il existe d'autres défaillances d'un appel de bas niveau, outre l'indicateur ignoré~: il réussit sur une adresse sans code, si bien qu'un appel destiné à un contrat inexistant paraît correct, sauf vérification. Il transmet par défaut presque tout le gas, de sorte que le contrat appelé peut le consommer, ou entrer en réentrance. Ses données de retour peuvent être volumineuses et coûteuses à copier. Un \texttt{delegatecall} exécute le code d'un autre sur le stockage du contrat appelant.

Dans une boucle, un seul bénéficiaire qui échoue (revert) fait échouer tout le paiement si l'on vérifie chaque appel, et une boucle qui s'allonge peut dépasser la limite de gas du bloc, si bien qu'elle ne se termine jamais. Si l'on ne vérifie pas, on enregistre des paiements qui n'ont pas eu lieu.

La réponse est le schéma de retrait, que recommande la documentation de Solidity~: on enregistre ce qui est dû à chacun, et chaque bénéficiaire retire lui-même son dû. Une défaillance ne touche alors que ce bénéficiaire. Le coût est une transaction supplémentaire par bénéficiaire, et quelqu'un doit la payer.

\iqwrong{«~C'est un paiement, et un paiement raté échoue (revert).~» Cela vaut pour \texttt{transfer}. Un \texttt{call} de bas niveau renvoie simplement false.}
\end{solution}

\begin{solution}{iq:iv3:smart-contracts-and-crypto-engineering:12}
Avec \texttt{byOrigin} à true, le piège prend les 5 ETH. Avec false, l'appel échoue (revert), et le portefeuille conserve ses 5 ETH. \texttt{tx.origin} est le compte qui a signé la transaction, et \texttt{msg.sender} est l'appelant immédiat. Le propriétaire a signé la transaction, si bien que \texttt{tx.origin} désigne toujours le propriétaire, même au bout d'une chaîne d'appels qui passe par le piège.

Ainsi, \texttt{tx.origin} répond à la question «~qui a déclenché tout cela~?~», alors qu'un agrément demande autre chose. Le propriétaire voulait appeler le piège et non le portefeuille. Une vérification sur \texttt{msg.sender} voit le piège comme appelant, et refuse. La documentation de Solidity indique explicitement de ne jamais utiliser \texttt{tx.origin} pour un agrément.

L'attaque suppose seulement que le propriétaire soit amené à signer une transaction quelconque qui atteint un contrat hostile, par exemple la réclamation d'un lot, ou un jeton doté d'un hook. Le propriétaire ne voit rien d'anormal dans ce qu'il a signé, à moins que le logiciel du portefeuille ne simule la transaction.

\texttt{tx.origin} pose aussi problème dans l'autre sens~: un portefeuille qui est lui-même un contrat, comme un portefeuille multisignature, ne signe jamais, si bien que \texttt{tx.origin} n'est pas ce portefeuille, et ne peut donc pas passer ce contrôle du tout.

\iqwrong{«~\texttt{tx.origin} est plus solide, car seule une personne réelle peut être l'origine.~» Être une personne au bout de la chaîne ne montre pas ce que cette personne voulait.}
\end{solution}

\begin{solution}{iq:iv3:smart-contracts-and-crypto-engineering:13}
Le montant maximal qu'il peut prendre est 160~: 100 avant le changement, puis 60 après, puisque le changement écrase simplement la valeur. Dans le contrat, il faut conditionner le changement à l'ancienne valeur, de sorte qu'une mise à jour périmée échoue, ou proposer des fonctions d'augmentation et de diminution. Côté client, il faut d'abord ramener l'allocation à 0, vérifier ce qui a été dépensé, et seulement ensuite fixer la nouvelle valeur.
\omcode{code/interviews-3/08-smart-contracts-and-crypto-engineering/python/iv_approve.py}{1}{19}{Une allocation, écrasée ou comparée.\protect\hfil\protect\break}\label{lst:iv3:cr:approve}

Le défaut tient à ce que \texttt{approve} ne sait rien de ce qui a déjà été dépensé. Le propriétaire voulait dire «~de 100 à 60 au total~», et le contrat entend «~60 à partir de maintenant~». Ainsi, un bénéficiaire qui voit la transaction en attente dépense d'abord les 100, puis dispose de 60 de plus. Le test confirme $n + m$ pour toute allocation, et montre qu'une comparaison donne $\max(n, m)$.

Avec une comparaison, le bénéficiaire peut toujours prendre ce qu'il pouvait prendre avant le changement, ce qui est inévitable puisqu'il y était autorisé. Il ne peut simplement pas tirer profit du changement. Passer par 0 produit le même effet seulement si le propriétaire vérifie entre les deux transactions ce qui a été dépensé, et réduit ou supprime en conséquence la nouvelle valeur. Fait à l'aveugle, cela donne encore 160 au bénéficiaire, comme le montre le test. Cela coûte deux transactions et une vérification, et c'est pourquoi la norme sur les jetons recommande aux clients de procéder ainsi. La norme précise aussi aux contrats de ne pas l'imposer, pour rester compatible avec les déploiements antérieurs.

Sur une chaîne publique, tout ce qui est visible pendant qu'une transaction est en attente peut être exploité en premier. Il ne faut pas écrire un changement d'état qui suppose que rien ne se passe entre la signature et l'inclusion. Un style plus récent repose sur des permits fondés sur une signature, avec un nonce et une échéance, de sorte que l'allocation s'épuise en une seule transaction~; les approbations illimitées accordées à des contrats comportent par ailleurs leur propre risque.

\iqwrong{«~Réduire une allocation est toujours sûr, car cela réduit le risque.~» Telle qu'elle est écrite, elle donne au bénéficiaire la possibilité de prélever sur les deux montants.}
\end{solution}

\begin{solution}{iq:iv3:smart-contracts-and-crypto-engineering:14}
1 ETH permet d'acheter 1992.01 dollars, contre un prix au comptant de 2000. Cela fait 0.40~\% de moins~: 0.3~\% de frais et environ 0.1~\% d'impact sur le prix. Le \Cref{lst:iv3:cr:amm} fait ce calcul sur des entiers. La division arrondit par défaut, si bien que l'acheteur reçoit un peu moins et le pool un peu plus, ce qui protège le pool. Le contrat et un modèle en Python concordent au wei près sur 101 cas.
\omcode{code/interviews-3/08-smart-contracts-and-crypto-engineering/solidity/Amm.sol}{4}{11}{Le montant de sortie d'un swap.\protect\hfil\protect\break}\label{lst:iv3:cr:amm}

Le pool empêche $x y$ de baisser. En prélevant les frais sur l'entrée, on obtient $a' = 0.997 a$, puis $(x + a')(y - b) = xy$ donne $b = a' y / (x + a')$. Le code multiplie par 997 et utilise 1000 au dénominateur, pour rester en entiers, et multiplie avant de diviser, pour conserver la précision.

Arrondir la sortie par défaut signifie que $x y$ ne peut que croître. Le sens compte~: s'il arrondissait par excès, un attaquant pourrait enchaîner un très grand nombre de petits swaps, en gagnant un wei à chaque fois. Le test vérifie sur 2000 pools aléatoires que l'invariant ne baisse jamais, et qu'un aller-retour ne peut pas rapporter de profit.

Pour vérifier que le contrat concorde avec le modèle, le test compile le contrat, l'exécute sur une EVM en mémoire avec des entrées aléatoires, et le compare exactement à une implémentation en Python qui reprend la même formule. Les deux partagent la formule~: ce contrôle détecte donc une erreur d'écriture ou un débordement, mais pas une mauvaise formule, et c'est à cela que servent les tests de propriétés.

Avec 18 décimales, le numérateur vaut ici environ $2 \times 10^{45}$, bien en dessous de $2^{256} \approx 10^{77}$, même si des jetons dont l'offre est très élevée peuvent s'en approcher~; un contrat réel peut donc avoir besoin d'une multiplication sur pleine largeur.

\iqwrong{«~Le prix vaut $y / x$, donc 1 ETH achète 2000.~» C'est le prix marginal, pour un échange infiniment petit, avant les frais.}
\end{solution}

\begin{solution}{iq:iv3:smart-contracts-and-crypto-engineering:15}
Le profit est d'environ 4881 dollars. Ces 9000 poussent le prix du pool à environ 100 fois sa valeur réelle. Le prêteur valorise 100 jetons à environ 9973, et prête 4986.5. La revente récupère 8995 des 9000. Avec un prix juste, l'attaquant perd 55. Le prix au comptant du pool est la faille~: sur une seule transaction, il peut prendre n'importe quelle valeur que l'on est prêt à payer.
\omcode{code/interviews-3/08-smart-contracts-and-crypto-engineering/python/iv_oracle.py}{1}{15}{Une transaction contre un prix au comptant.\protect\hfil\protect\break}\label{lst:iv3:cr:oracle}

Toutes les étapes ont lieu dans une même transaction~: emprunter, acheter, déposer le collatéral, emprunter contre lui, revendre, rembourser. L'attaquant repart avec le prêt et abandonne un collatéral valant 100, tandis que le prêteur supporte une mauvaise créance. Comme le premier prêt est remboursé dans la même transaction, l'attaquant n'a eu besoin d'aucun capital~; si une étape échoue, l'ensemble est annulé et seul le gas est perdu. L'attaque ne coûte donc rien à tenter.

Le flash loan n'est pas la cause profonde, puisqu'il supprime seulement le besoin de capital. La cause profonde est de traiter comme vérité un prix qu'il est facile de déplacer.

Défenses~: utiliser un prix qu'une seule transaction ne peut pas déplacer. Il peut s'agir d'une moyenne dans le temps prise sur le pool, d'un flux de prix agrégeant plusieurs plateformes de négociation, ou d'une limite sur l'ampleur d'un mouvement. La taille compte aussi~: sur un pool profond, la même attaque perd de l'argent, comme le montre le test. Un prêteur peut aussi limiter le montant empruntable contre un collatéral peu liquide, et retarder un emprunt d'un bloc, ce qui a cependant un coût.

Le modèle suppose un prêt sans coût et sans gas~; il ne donne qu'une esquisse du mécanisme sur des chiffres locaux, sans aucun contrat impliqué.

\iqwrong{«~Elle n'affecte que la transaction de l'attaquant, puisque le prix revient aussitôt.~» Le prêteur a lu le prix au moment où il était faux, et ce prêt subsiste.}
\end{solution}

\begin{solution}{iq:iv3:smart-contracts-and-crypto-engineering:16}
La moyenne vaut 1.99, contre un prix au comptant de 100 pendant le pic. Une transaction isolée déplace peu la somme, puisqu'elle ne dure que peu de temps. Pour déplacer la moyenne, un attaquant doit maintenir le prix à l'écart pendant une grande partie de la fenêtre, pendant que d'autres échangent à ce prix. Une fenêtre plus longue est plus difficile à déplacer, mais suit le marché plus lentement.
\omcode{code/interviews-3/08-smart-contracts-and-crypto-engineering/python/iv_twap.py}{1}{17}{Un accumulator de prix.\protect\hfil\protect\break}\label{lst:iv3:cr:twap}

L'accumulator coûte peu. Il stocke un seul total cumulé, mis à jour au premier accès à un bloc, si bien qu'un contrat lit la somme à deux instants et la divise par le temps écoulé. Rien n'est stocké à chaque seconde. Le test compare ce résultat à la moyenne de référence calculée seconde par seconde, pour 300 trajectoires aléatoires.

Pour tenir le prix, un attaquant doit le maintenir élevé d'un bloc à l'autre, de sorte qu'il paie les arbitragistes à chaque bloc, et ce coût croît avec la profondeur et avec le temps. Et au sein d'un même bloc, une lecture au début ne tient pas compte des modifications faites dans ce bloc, ce qui est un élément clé de la conception. C'est une défense, et non une garantie~: un producteur de plusieurs blocs consécutifs peut éviter l'arbitrage intermédiaire, et un pool très mince reste peu coûteux à faire monter.

La fenêtre arbitre entre sécurité et retard. Lors d'un krach brutal, une moyenne sur une heure retarde sur le marché, si bien qu'un prêteur valoriserait le collatéral trop haut justement au moment où cela compte, et cela peut aussi être exploité.

La moyenne arithmétique donne trop de poids à un pic~: dans notre exemple, un pic multiplié par 100 pendant 1~\% du temps suffit à la faire doubler. Une moyenne du logarithme du prix, autrement dit une moyenne géométrique, compte de façon égale une hausse et une baisse de même ampleur, et donne ici $100^{0.01} = 1.05$. Une conception peut donc stocker celle-ci à la place.

\iqwrong{«~Les moyennes ne peuvent pas être manipulées.~» Elles coûtent seulement plus cher à manipuler, et le surcoût dépend de la profondeur du pool et de la fenêtre.}
\end{solution}

\begin{solution}{iq:iv3:smart-contracts-and-crypto-engineering:17}
\texttt{trusting} renvoie 200\,000 dans les deux premiers cas, périmé ou non. Après $-1$, il renvoie $2^{256} - 1$, car la conversion d'un négatif en non signé s'enroule. \texttt{careful} renvoie 200\,000 à $T + 60$, et échoue (revert) à $T + 7200$ et sur $-1$. Il faut aussi vérifier les décimales, la complétude du round et l'état propre du flux.

La péremption des données est la défaillance la plus fréquente. Les flux publient lorsque le prix varie d'un seuil fixé ou qu'un délai donné s'est écoulé, si bien que la bonne limite dépend du heartbeat de ce flux. Une limite nettement plus longue que le heartbeat ajoute du risque, et une limite plus courte échoue (revert) pendant les périodes calmes normales.

Autres vérifications~: les décimales du flux, qui diffèrent d'un flux à l'autre et d'un jeton à l'autre, puisqu'une erreur à ce niveau fausse le prix d'une puissance de dix. Les bornes minimale et maximale d'un flux, auxquelles un agrégateur peut cesser de suivre le marché et se contenter de publier la borne. Sur un layer 2, il faut vérifier que le séquenceur est en service, car personne ne pouvait publier de mises à jour pendant son arrêt, et les prix peuvent être anciens à son retour.

Ce qu'il faut faire après un refus demande réflexion, et le revert n'est pas toujours la bonne réponse. Un prêteur qui annule sur un prix erroné a aussi bloqué les remboursements et les liquidations, justement quand on en a besoin. Il ne faut donc suspendre que les actions risquées, comme les nouveaux emprunts, et laisser les utilisateurs réduire leur risque. Une seconde source peut servir de repli, mais elle doit alors faire l'objet du même examen.

\iqwrong{«~Le flux est un contrat de confiance, donc sa réponse est bonne.~» Il publie ce qu'il publie, et c'est au code appelant de vérifier que la donnée est fraîche et exploitable.}
\end{solution}

\begin{solution}{iq:iv3:smart-contracts-and-crypto-engineering:18}
Avec une limite de 1~\%, la plus grosse première transaction vaut environ 10\,141 dollars, et le profit environ 140 dollars. Le profit croît avec la limite~: environ 70 à 0.5~\% et 420 à 3~\%. Il s'agit d'une attaque sandwich, une forme de MEV. Un utilisateur peut fixer une limite serrée, soumettre sa transaction en privé ou fractionner l'opération. Une limite de 0 échoue (revert) sur presque toute autre transaction.
\omcode{code/interviews-3/08-smart-contracts-and-crypto-engineering/python/iv_sandwich.py}{1}{29}{Une transaction avant et après la victime.\protect\hfil\protect\break}\label{lst:iv3:cr:sandwich}

Le trader achète en premier, ce qui fait monter le prix que la victime paie ensuite. La transaction propre de la victime le fait monter encore, et le trader revend dans ce mouvement. La limite de la victime fixe jusqu'où le trader peut aller~: la dichotomie trouve la première transaction qui laisse la victime exactement à son minimum, soit 9.77 jetons contre 9.87 cotés. Une recherche sur grille confirme que cette taille donne le profit maximal dans cette plage, et le test montre que, sans victime, les deux jambes ne font que payer des frais.

La limite est donc le budget de l'attaquant. Elle correspond à peu près à ce que l'utilisateur cède, si bien qu'une limite lâche est une invitation.

Une limite à 0 protège totalement contre cette attaque, mais elle fait qu'une transaction ordinaire passée en premier fait échouer (revert) la vôtre, et une transaction annulée paie quand même le gas. Le choix oppose donc l'échec à la perte. La soumission privée par un constructeur de blocs masque la transaction aux regards publics, au prix de la confiance accordée à ce canal. Fractionner une transaction, ou recourir à une plateforme d'enchères, réduit ce qui peut être prélevé.

Pour un teneur de marché, la même logique s'applique. Vos propres swaps en attente sont visibles, si bien que la taille et les limites font partie de l'exécution, et les profits calculés ici ignorent le gas ainsi que tout pot-de-vin nécessaire pour remporter la position, qui en représente souvent l'essentiel.

\iqwrong{«~Réglez le slippage à un niveau élevé pour que la transaction passe toujours.~» Elle passe, et cette marge revient à celui qui surveille.}
\end{solution}

\begin{solution}{iq:iv3:smart-contracts-and-crypto-engineering:19}
La forme ABI occupe 68 octets et 632 gas de données. Une forme compactée en occupe 29 et 440 gas. Les rollups s'en soucient, car les données publiées sur le layer 1 constituent depuis longtemps la part principale de leur coût. Les frais d'un layer 2 sont son propre coût d'exécution, auquel s'ajoute une part du coût sur le layer 1 de la publication de ces données.
\enlargethispage{-3\baselineskip}%
\omcode{code/interviews-3/08-smart-contracts-and-crypto-engineering/python/iv_calldata.py}{1}{13}{Ce que coûtent les octets d'une transaction.\protect\hfil\protect\break}\label{lst:iv3:cr:calldata}

Dans la forme ABI, chaque argument occupe 32 octets~: les 4 du sélecteur, puis 12 zéros et une adresse de 20, puis un nombre presque entièrement nul. Les zéros coûtent moins cher, ce qui atténue l'écart~: le nombre d'octets baisse de 57~\%, et le gas de seulement 30~\%. Le reste tient à l'adresse, dont les octets non nuls coûtent plein tarif. La forme compactée exige son propre décodeur dans le contrat, et elle peut omettre la longueur du montant, déductible puisque ce montant vient en dernier. Un vrai layer 2 peut aussi compresser avant de publier, ce qui modifie encore les calculs.

Les règles changent, et celle-ci est à revérifier. La règle 16-et-4 est l'ancienne~; depuis, une borne plancher a été ajoutée pour les transactions chargées en données, à 10 gas par jeton, un octet non nul valant 4 jetons. Notre transfert pèse 158 jetons, donc sa borne plancher vaut 1580 gas, soit moins qu'une écriture dans un emplacement, et cette borne ne joue pas. Par ailleurs, les blobs offrent aux rollups un canal moins cher, avec son propre marché des frais, pour des données qui ne restent pas disponibles indéfiniment et que les contrats ne peuvent pas lire. Le poids des données d'appel sur un layer 2 donné dépend donc de la manière dont il publie, et il faut consulter sa formule de frais en vigueur.

Ce qui change dans votre code~: sur un layer 2, l'exécution est souvent bon marché et les données coûteuses, si bien que les réflexes de stockage contre données d'appel du layer 1 peuvent s'inverser, et les encodages courts sont rentables. Les numéros de bloc et les horodatages peuvent avoir un sens différent, et le séquenceur fixe l'ordre, si bien que la finalité et la MEV se présentent autrement. Une cotation de frais doit aussi lire le coût actuel sur le layer 1, et peut être périmée au moment de l'inclusion.

\iqwrong{«~Sur un layer 2, comme sur le layer 1, le gas par opcode détermine le coût.~» Les frais d'un simple transfert y sont peut-être surtout constitués de données.}
\end{solution}

\begin{solution}{iq:iv3:smart-contracts-and-crypto-engineering:20}
\Cref{lst:iv3:cr:deposits} reconstruit la meilleure chaîne à partir de tous les blocs vus. Un dépôt est crédité une fois qu'il se trouve à $N$ blocs de profondeur, identifié par son propre identifiant et non par bloc, de sorte qu'une transaction qui réapparaît sur la nouvelle branche ne peut pas être créditée deux fois. Une réorganisation plus profonde que $N$ retire un dépôt déjà crédité, et cela déclenche une alarme pour une personne.
\omcode{code/interviews-3/08-smart-contracts-and-crypto-engineering/python/iv_deposits.py}{1}{32}{Dépôts et réorganisations.\protect\hfil\protect\break}\label{lst:iv3:cr:deposits}

C'est la clé qui rend l'opération idempotente. Un hachage de bloc ne suffit pas, car le même dépôt figurant dans un autre bloc a un autre hachage. Dans un système réel, l'identifiant est le hachage de la transaction, complété par l'indice du log, puisqu'une même transaction peut contenir plusieurs dépôts. L'état du processeur et le crédit doivent être écrits dans le registre en une seule transaction, et son redémarrage reprend à partir de la chaîne stockée.

Une réorganisation peu profonde, en deçà de $N$, n'exige aucune action, puisque rien n'a encore été crédité. C'est pourquoi $N$ se choisit en fonction des profondeurs réellement observées sur cette chaîne. Au-delà de $N$, le dépôt peut avoir été perdu définitivement, et l'utilisateur peut déjà l'avoir échangé ou retiré. Le processeur retire le crédit et déclenche l'alarme, et seule une personne ou une politique peut décider de la suite. C'est pourquoi les dépôts importants ou peu sécurisés doivent attendre plus longtemps, et que le retrait doit intervenir après le trading.

Le test couvre des cas fixes~: trop peu de confirmations, un seul crédit, une réorganisation profonde avec alarme, et une transaction minée de nouveau sur la nouvelle branche. Il construit aussi 300 arbres aléatoires dont les bifurcations sont plus courtes que $N$, et vérifie que chaque dépôt à la profondeur $N$ ou au-delà sur la meilleure chaîne finale a été crédité.

Le modèle retient le bloc le plus haut comme meilleur, alors qu'un client réel suit la règle de bifurcation de la chaîne. Il parcourt toute la chaîne à chaque bloc par souci de clarté, alors qu'un client réel ne garderait qu'une fenêtre. Il suppose aussi avoir vu chaque bloc~: après une interruption, un système réel doit combler le retard et traiter un trou comme une erreur, faute de quoi un dépôt manqué ne sera jamais crédité.

\iqwrong{«~Écoutez l'événement de dépôt et créditez-le.~» Un événement issu d'un bloc qui disparaît ensuite est un crédit sans contrepartie.}
\end{solution}

\begin{solution}{iq:iv3:smart-contracts-and-crypto-engineering:21}
Une nouvelle tentative doit conserver son nonce, car un nouveau nonce donne une seconde transaction, qui pourrait elle aussi être minée, et l'utilisateur serait payé deux fois. Avec le même nonce, une seule peut jamais aboutir. Un nonce bloqué bloque toutes les demandes qui le suivent. Relever les frais en renvoyant ce nonce à un prix supérieur. Enregistrer l'association entre la demande et le nonce avant la diffusion.
\omcode{code/interviews-3/08-smart-contracts-and-crypto-engineering/python/iv_withdraw.py}{1}{36}{Retraits et relèvements de frais.\protect\hfil\protect\break}\label{lst:iv3:cr:withdraw}

Chaque demande est donc liée une fois pour toutes à un nonce. Le passage aléatoire du test relève les frais au hasard, avec des frais de marché mouvants, et vérifie que chaque demande est payée exactement une fois. Il montre aussi le blocage~: un nonce bas bon marché retient une transaction à frais élevés qui le suit, jusqu'à la correction. La copie défectueuse, qui donne un nouveau nonce à la nouvelle tentative, paie deux fois.

Enregistrer d'abord compte en cas de plantage. Si vous diffusez puis plantez avant d'avoir enregistré, un redémarrage peut ignorer ce qui est parti. Notez donc la demande, le nonce et les frais, puis diffusez, et au redémarrage rediffusez tout ce qui n'est pas confirmé. Rediffuser est sans danger, car le même nonce avec le même contenu ne peut être miné qu'une fois.

Avec un seul compte, une transaction bloquée retarde tout le monde~; les opérateurs peuvent donc utiliser plusieurs comptes chauds en parallèle. Une politique de frais élevés, avec des relèvements à intervalles fixes, fait avancer la file. En pratique, le remplacement doit dépasser l'ancien montant de frais d'une marge qui dépend du nœud, et des frais très bas peuvent être purement et simplement retirés du pool, si bien qu'une boucle de rediffusion est nécessaire.

La gestion des confirmations rejoint sur ce point la logique des dépôts~: un retrait miné peut être sorti par réorganisation, puis revenir, ou devoir être renvoyé avec le même nonce. Ne considérez un retrait comme définitif qu'après une profondeur suffisante. Et une cible contractuelle peut faire échouer la transaction, en consommant le nonce sans payer l'utilisateur~: lisez le reçu, et ne vous contentez pas de vérifier qu'elle a été minée.

\iqwrong{«~Elle semble bloquée, alors renvoyez-la avec le nonce suivant.~» C'est ainsi que des plateformes d'échange paient deux fois.}
\end{solution}

\begin{solution}{iq:iv3:smart-contracts-and-crypto-engineering:22}
L'écriture part dans l'emplacement 0 du proxy, qui est son \texttt{impl}, puisque l'implémentation considère que l'emplacement 0 est \texttt{x}. Ainsi \texttt{impl()} renvoie alors \texttt{0xdead}. Un appel de \texttt{x()} à travers le proxy délègue désormais à une adresse sans code, qui réussit et ne renvoie rien. Placer l'adresse de l'implémentation dans un emplacement choisi par hachage, et n'ajouter des variables qu'à la mise à niveau.

\texttt{delegatecall} exécute le code d'un autre contrat en utilisant le stockage, le solde et \texttt{msg.sender} de l'appelant. Le code retrouve ses variables par numéro d'emplacement, sans rien savoir de ce que le proxy y conserve. Ainsi \texttt{x = v} signifie «~écrire l'emplacement 0~», quel que soit le lieu d'exécution du code. C'est là le défaut~: deux contrats partagent un même stockage, mais ont chacun une idée différente de sa disposition.

L'exécution locale le confirme. Le proxy finit avec son adresse remplacée, sans qu'aucune erreur ne soit signalée à aucune étape, et plus grave encore, chaque appel ultérieur semble fonctionner. Une cible choisie par un attaquant lui donnerait accès à tous les fonds du proxy.

La réparation standard utilise un emplacement dérivé d'un hachage, qu'aucune variable attribuée par le compilateur n'atteindra. Notre \texttt{SafeProxy} procède ainsi~: le même appel écrit alors \texttt{x} dans l'emplacement 0, tandis que l'adresse de l'implémentation reste en place. Un tel standard existe pour que les outils puissent lire cet emplacement également.

Une implémentation mise à niveau doit conserver la disposition de l'ancienne~: pas de réordonnancement, pas de types modifiés, et des ajouts uniquement en fin. Les constructeurs ne s'exécutent pas dans le contexte du proxy~: l'état est donc fixé par un initialiseur, qui ne doit pouvoir être appelé qu'une seule fois, car un initialiseur ouvert permet une prise de contrôle. Une fonction du proxy peut aussi partager un sélecteur avec une fonction de l'implémentation, et la fonction de mise à niveau exige un contrôle d'accès.

\iqwrong{«~L'implémentation a son propre stockage, donc le proxy est sûr.~» Sous \texttt{delegatecall}, le stockage de l'implémentation n'est jamais touché.}
\end{solution}

\begin{solution}{iq:iv3:smart-contracts-and-crypto-engineering:23}
\texttt{Reader} lit 0.5, contre 1.0 avant et après. Dans le rappel, les ETH ont déjà quitté le contrat, alors que \texttt{shares} n'a pas encore baissé~: \texttt{price} divise donc le plus petit solde par l'ancienne offre. C'est une réentrance, puisque le pool est réappelé alors qu'il n'est pas en phase avec son état. Un prêteur qui utiliserait ce prix sous-évaluerait la garantie, ou la surévaluerait dans une conception symétrique. Corriger en mettant à jour l'état en premier, ou en faisant refuser la lecture par une vue.

Aucun état ne change dans le rappel, puisque \texttt{price} est une vue. C'est pourquoi on parle de réentrance en lecture seule, et pourquoi elle échappe facilement~: un verrou sur \texttt{leave} et \texttt{join} empêche de rentrer de nouveau dans ces fonctions, mais il n'a aucun effet sur une vue qui ne vérifie pas le verrou. Le pool lui-même ne perd rien. La victime est un autre contrat, qui se fie au prix du pool.

Que perdrait un prêteur~? S'il lit un prix bas, les positions saines paraissent sous-garanties, et elles peuvent être liquidées à un faux prix pendant ce même rappel. Dans une conception symétrique, où le prix grimpe en plein appel, un attaquant peut emprunter trop. Le coût de l'attaquant se limite au gas et au capital temporaire, qui peut provenir d'un flash loan.

Notre correction fait les deux choses~: \texttt{shares} baisse avant l'appel, si bien que l'état reste juste à chaque instant, et \texttt{price} échoue (revert) tant que le verrou est posé. Le \texttt{SafeReader} du test est refusé, et l'exécution se termine tout de même avec un prix de 1.

Un consommateur dispose lui aussi d'une défense propre~: vérifier le verrou de l'autre contrat, s'il en expose un, avant de faire confiance à sa vue. Les revues doivent traiter toute fonction de vue utilisée par d'autres comme faisant partie de la sûreté du contrat, et non comme une simple lecture sans conséquence.

\iqwrong{«~C'est une fonction de vue, elle ne peut donc pas faire partie d'une attaque.~» Elle ne modifie rien par elle-même. Le mal vient de ce qu'on agit sur ce qu'elle a affirmé.}
\end{solution}

\begin{solution}{iq:iv3:smart-contracts-and-crypto-engineering:24}
Les trois soumissions réussissent~: la première paie, la rejouée paie encore, et l'autre déploiement paie aussi. Le message signé ne contient que le destinataire et le montant. Il n'a pas de nonce, il peut donc être rejoué. Il ne contient ni l'adresse du contrat ni l'identifiant de chaîne, il fonctionne donc ailleurs. Un utilisateur peut payer sans gas, puisque le relayeur soumet et paie, et le recouvrement des coûts se fait autrement.
\omcode{code/interviews-3/08-smart-contracts-and-crypto-engineering/solidity/Relay.sol}{20}{46}{La correction.}\label{lst:iv3:cr:saferelay}

Une signature prouve que quelqu'un a approuvé ces octets exacts. Elle ne dit rien du nombre de fois ni du lieu, sauf si ces octets le précisent. Chaque champ bloque donc un type de rejeu. Le nonce empêche les répétitions sur ce contrat, et l'adresse du contrat empêche son usage sur un autre déploiement. L'identifiant de chaîne empêche l'usage sur une autre chaîne, par exemple après une bifurcation, ou une copie à la même adresse ailleurs. Le signataire figure lui aussi, de sorte que \texttt{ecrecover} doit renvoyer cette adresse. Les normes de données typées ajoutent un séparateur de domaine qui contient exactement ces champs, et elles laissent la protection contre le rejeu à l'application, d'où le nonce.

Nous avons exécuté les deux versions sur l'EVM local, avec une vraie signature issue d'une clé de test. Le premier relais a payé Mallory deux fois, et aussi sur la seconde copie. La version corrigée a payé une fois, et a refusé à la fois le rejeu et l'autre déploiement, ce qui prouve que l'empreinte lie bien ces éléments.

Quelques précisions. \texttt{ecrecover} renvoie zéro pour une mauvaise signature, si bien qu'un zéro non vérifié peut correspondre à un compte non initialisé, et la correction le refuse. Les signatures sont malléables~: il existe un autre couple $(r, s)$ valide pour le même message, il ne faut donc pas utiliser la signature elle-même comme clé anti-rejeu, car un nonce le fait correctement. Ajouter aussi une échéance, car une signature inutilisée resterait valable indéfiniment.

Le relayeur paie le gas, et il se rembourse sur le transfert, ou par un sponsor. C'est ce que signifie «~pas de gas pour l'utilisateur~»~: quelqu'un paie quand même.

\iqwrong{«~La signature est valide, donc le paiement est autorisé.~» Valide pour ce message, et ce message ne dit pas combien de fois.}
\end{solution}

\begin{solution}{iq:iv3:smart-contracts-and-crypto-engineering:25}
Avec un schéma 2-sur-3, la probabilité de vol vaut $3p^2 - 2p^3 \approx 0.03~\%$, contre 1~\% pour une seule clé, et les fonds survivent à la perte d'un détenteur. Avec $t - 1$ parts, chaque secret reste possible, puisqu'un polynôme de degré $t - 1$ peut toujours passer par n'importe quelle valeur en $0$. Le \Cref{lst:iv3:cr:shamir} est le code.
\omcode{code/interviews-3/08-smart-contracts-and-crypto-engineering/python/iv_shamir.py}{1}{28}{Partage de secret.}\label{lst:iv3:cr:shamir}

Le secret est la valeur du polynôme en 0, les autres coefficients étant aléatoires. Chaque part est un point de ce polynôme. $t$ points déterminent un polynôme de degré $t - 1$, par interpolation. Avec moins de points, pour chaque secret candidat il existe exactement un polynôme passant par les points connus, si bien que dans un corps fini tous les secrets sont également probables. Le système n'est donc pas seulement difficile à casser, il est impossible à casser, quelle que soit la puissance de calcul. Le test énumère ce cas pour $t = 2$ sur le corps de taille 257. Le même test retrouve le secret à partir de chaque sous-ensemble de 3 parts parmi 5. Le corps doit être plus grand que le secret~: une clé de 32 octets exige un nombre premier supérieur à $2^{256}$, si bien que le programme utilise $2^{521} - 1$, et \texttt{split} refuse un secret qui ne tient pas, car le réduire perdrait la clé sans le signaler.

Le calcul du vol est une énumération~: 3 façons de compromettre deux détenteurs, chacune de probabilité $p^2(1-p)$, et une façon de les compromettre tous les trois, de probabilité $p^3$. L'indépendance est l'hypothèse à attaquer, car les détenteurs peuvent partager un logiciel, un bâtiment ou un processus.

La faiblesse de Shamir tient à ce que la clé est reconstituée en un seul endroit au moment de l'usage, et qu'à cet instant une seule machine la détient. Un contrat multisignature utilise des clés distinctes, et chacune signe, si bien qu'il n'existe jamais de clé unique~; la règle est visible sur la chaîne, au prix de davantage de gas, et seulement là où la chaîne le permet. La signature à seuil (MPC) fait produire conjointement une signature ordinaire aux parties, sans que la clé n'existe jamais en un seul endroit. Cela exige un protocole bien plus complexe, dont l'implémentation constitue en soi un risque, et c'est là qu'un système de conservation réel doit s'appuyer sur du code audité. Quel que soit votre choix, ce code est destiné à l'enseignement~: \texttt{secrets} fournit l'aléa, mais les canaux auxiliaires et le stockage des parts ne sont pas traités.

\iqwrong{«~Coupez la clé de 32 octets en trois morceaux, et donnez-en deux à chaque détenteur.~» Alors un détenteur connaît une partie de la clé, et chaque vol réduit la recherche.}
\end{solution}

\begin{solution}{iq:iv3:smart-contracts-and-crypto-engineering:26}
\Cref{lst:iv3:cr:merkle} est une réponse, avec des preuves d'au plus 5 hachages pour 23 comptes. Les préfixes empêchent qu'une feuille soit prise pour un nœud interne. Une preuve établit que ce solde figure dans l'arbre derrière la racine publiée. Elle n'établit pas que la bourse détient les actifs, ni que tous les passifs sont dans l'arbre. La solvabilité exige aussi une preuve des actifs, et un total.
\omcode{code/interviews-3/08-smart-contracts-and-crypto-engineering/python/iv_merkle.py}{1}{37}{Un arbre de soldes.\protect\hfil\protect\break}\label{lst:iv3:cr:merkle}

Une preuve comprend un frère par niveau, si bien que sa longueur vaut au plus $\lceil \log_2 23 \rceil = 5$. L'utilisateur recalcule en remontant à partir de sa propre feuille, et compare le résultat à la racine publiée. La bourse ne peut pas modifier un solde sans modifier la racine, et l'utilisateur n'a besoin que des hachages.

Sans préfixes, l'entrée d'un nœud interne, formée de deux hachages de 32 octets, est aussi une entrée légale de feuille de 64 octets~; on pourrait donc présenter un nœud interne comme s'il était une feuille. Des préfixes distincts rendent ces deux types de hachage disjoints.

Ce que la preuve ne démontre pas compte davantage. Seuls les utilisateurs qui vérifient sont protégés, car la bourse pourrait omettre ceux qu'elle croit peu susceptibles de contrôler. L'arbre ne somme pas les valeurs~: il faut donc ajouter à chaque nœud un total cumulé, de sorte que le total d'un parent soit la somme de ses enfants, et que la racine engage le total des passifs. Un compte au solde négatif pourrait réduire le total~; une preuve doit donc établir que tous les soldes sont positifs ou nuls, et les preuves à divulgation nulle de connaissance peuvent le faire tout en préservant la vie privée. Enfin, un instantané ne représente qu'un seul moment~: des fonds pourraient être empruntés pour ce seul moment. La vie privée est aussi en jeu, car une feuille contenant un nom de compte et un solde en clair divulgue ces informations~; on utilisera donc un hachage salé.

La solvabilité exige aussi les actifs~: une signature émanant des portefeuilles, ou un mouvement sur la chaîne, prouvant le contrôle. Les prêts consentis sur ces actifs ne sont pas non plus montrés. Un audit, répété souvent, est ce qui relie ces éléments entre eux. Le test vérifie les 23 preuves, rejette chacune où le solde a été modifié, et couvre le cas d'un arbre à une seule feuille. La vue propre à la chaîne sur ces actifs ne figure pas dans ce modèle.

\iqwrong{«~Une preuve de réserves signifie que la bourse est solvable.~» Des actifs sans passifs ne prouvent rien, et cet arbre seul ne montre qu'une partie des passifs.}
\end{solution}

\begin{solution}{iq:iv3:smart-contracts-and-crypto-engineering:27}
Le prix de liquidation vaut $(2000 - 200) / (1 - 0.005) = 1809.05$. Un prix de marquage résiste à une transaction isolée anormale ou à un carnet peu profond, qui pourraient autrement déclencher des clôtures. Si la clôture s'exécute à 1790, les fonds propres valent $-10$, et le fonds d'assurance couvre cet écart. Si le fonds s'épuise, les traders bénéficiaires de l'autre côté peuvent subir une réduction de position, par désendettement automatique, ou une part de la perte.
\enlargethispage{-4\baselineskip}%
\omcode{code/interviews-3/08-smart-contracts-and-crypto-engineering/python/iv_liquidation.py}{1}{22}{Liquidation.\protect\hfil\protect\break}\label{lst:iv3:cr:liquidation}

Les fonds propres valent la marge plus le profit latent~: $m + (P - E) q$. En posant ce résultat égal à $r P q$ et en résolvant, on obtient $P = (E q - m) / (q (1 - r))$. Le test le retrouve aussi par dichotomie sur \texttt{must\_close} pour 500 positions aléatoires, et les deux résultats concordent. Avec une marge initiale de 10~\%, une baisse d'environ 9.5~\% entraîne la clôture de la position.

Le prix de marquage provient en général d'un indice regroupant plusieurs plateformes au comptant, auquel s'ajoute une base amortie. Si la liquidation reposait sur la dernière transaction, quelqu'un pourrait pousser un instant le carnet peu profond d'une plateforme, clôturer les positions des autres et prendre la contrepartie. Des prix de marquage équitables suppriment l'essentiel de ce jeu, mais l'indice lui-même doit résister à la défaillance d'une plateforme.

La séquence est la suivante~: le moteur clôture sur le marché au moyen d'un ordre agressif, idéalement bien avant que les fonds propres n'atteignent 0, et toute marge restante revient au trader ou au fonds. La marge de maintenance constitue le coussin contre une baisse pendant la clôture. En cas d'écart, ou de carnet peu profond, l'exécution se fait en dessous du seuil de faillite, et quelqu'un doit absorber le déficit, puisque le gagnant de l'autre côté est créancier de cette somme. C'est pourquoi les bourses alimentent un fonds à partir des excédents de liquidation.

Une bourse doit faire fonctionner ce mécanisme sous la charge la plus lourde qu'elle rencontrera, lorsque de nombreuses positions sont clôturées en même temps et que chaque clôture pousse le prix plus loin. Les moteurs clôturent donc par étapes, limitent la vitesse de vente, et peuvent relever les marges. Ces clôtures forment la cascade classique. Un prix de marquage erroné ou un flux périmé peut causer de graves dégâts ici, et doit donc être traité comme les contrôles de risque du chemin critique. Ce qui précède vaut pour la marge isolée~; la marge croisée compense les positions entre elles.

\iqwrong{«~La liquidation a lieu lorsque les fonds propres atteignent zéro, donc la bourse ne perd jamais.~» Au moment où l'ordre s'exécute, le prix peut être bien au-delà de ce seuil.}
\end{solution}

\begin{solution}{iq:iv3:smart-contracts-and-crypto-engineering:28}
La taille optimale vaut environ 6979 dollars, pour un profit d'environ 24.28 avant les frais de gas et de plateforme. En dessous de $2000 / 0.997 = 2006.02$, il n'y a aucune transaction, car les frais absorbent l'écart. La formule ne tient compte ni du calendrier, ni du risque d'échec ou d'être devancé, ni du gas, ni de ce que l'on détient pendant l'attente.
\enlargethispage{-2\baselineskip}%
\omcode{code/interviews-3/08-smart-contracts-and-crypto-engineering/python/iv_arb_size.py}{1}{14}{La taille d'un arbitrage.\protect\hfil\protect\break}\label{lst:iv3:cr:arb}

Avec $k = 1 - f$, les jetons reçus pour $d$ dollars valent $k d x / (y + k d)$. Le profit est le prix extérieur $S$ multiplié par ce montant, diminué de $d$. En annulant la dérivée, on obtient $(y + k d)^2 = k x y S$, d'où $d^* = (\sqrt{k x y S} - y) / k$. Le test compare ce résultat à une recherche sur grille, et confirme le seuil, car $d^*$ n'est positif que si $k S x > y$. Après la transaction, le prix marginal du pool, net des frais, est égal au prix extérieur.

Le gas est un coût fixe, il ne modifie donc pas $d^*$. Il décide seulement si l'on trade~: ici, un coût de gas supérieur à 24.28 signifie qu'on ne trade pas. Le test couvre ce cas. Les frais d'une plateforme centralisée abaisseraient en pratique $S$, et déplaceraient $d^*$.

Les risques sont l'enjeu même de l'entretien. Les deux jambes ne peuvent pas être simultanées, car un bloc arrive toutes les quelques secondes tandis que la plateforme négocie en continu~; le prix extérieur peut donc bouger, et l'on pourrait vendre d'abord et acheter ensuite, ou l'inverse. D'autres voient la même opportunité, et l'enchère de priorité peut absorber l'essentiel du profit sous forme de pot-de-vin. Une transaction échouée coûte tout de même du gas. Il faut détenir un inventaire des deux côtés, et le déplacer prend du temps, avec un risque de chaîne pendant le transfert. Un système réel doit donc valoriser tous ces éléments et fixer une limite pour chacun, et il peut couvrir l'écart jusqu'à ce que les deux jambes soient exécutées. Nous n'avons fait tourner que le modèle Python ci-dessus.

\iqwrong{«~Achetez jusqu'à ce que le prix du pool soit égal au prix extérieur.~» Cela ignore les frais. L'optimum laisse un écart de la taille des frais.}
\end{solution}

\begin{solution}{iq:iv3:smart-contracts-and-crypto-engineering:29}
Tenir un grand livre interne en partie double. Un ordre ouvert bloque des fonds prélevés sur le solde libre de l'utilisateur, de sorte qu'ils ne puissent pas être dépensés deux fois. Le \Cref{lst:iv3:cr:holds} en est le cœur. Le moteur gère l'appariement, le grand livre gère l'argent, et les deux se relient par un flux séquencé. Conserver peu de fonds dans les portefeuilles chauds. Il ne doit jamais être possible de créer de la monnaie, de dépenser des fonds bloqués, ni de payer deux fois un retrait.
\omcode{code/interviews-3/08-smart-contracts-and-crypto-engineering/python/iv_holds.py}{1}{32}{Soldes libres et bloqués.\protect\hfil\protect\break}\label{lst:iv3:cr:holds}

Un ordre passe d'abord par un blocage dans le grand livre, puis est transmis au moteur d'appariement. Un blocage doit être effectué avant que le moteur ne le voie. Les transactions du moteur reviennent sur un flux séquencé, chacune réglant le blocage du payeur vers le solde libre du bénéficiaire, et une annulation libère le reliquat. Le test soumet à ce système des opérations aléatoires, et vérifie que la somme des soldes libres, bloqués et retirés reste constante, et qu'aucun solde ne devient négatif. Pour des raisons de rapidité, un moteur conserve en général ses propres soldes en mémoire, dans le chemin séquencé, le grand livre suivant le même journal, de sorte que les deux soient rapprochés et puissent être reconstruits.

Les dépôts passent par un surveillant de chaîne vers le grand livre, et ne sont crédités qu'une fois suffisamment profonds pour qu'une réorganisation soit peu susceptible de les annuler, une profondeur fixée pour chaque chaîne. Les retraits passent par des contrôles de risque et des limites, sont bloqués dans le grand livre, puis sont signés par un service distinct, sans aucun chemin direct entre un serveur web et une clé. Chaque retrait est lié, pour toujours, par son identifiant de requête à un seul nonce, de sorte qu'une nouvelle tentative ne puisse que relever les frais, et ne puisse pas payer deux fois.

Les portefeuilles chauds ne contiennent que ce qu'exigent les sorties nettes récentes, et un transfert vers le stockage tiède ou froid suit un calendrier. Supposons que les retraits nets atteignent un pic de 2~\% des fonds en un jour, et que le réapprovisionnement depuis le stockage froid prenne 6 heures~: la réserve chaude doit alors détenir environ 0.5~\%, plus une marge. Il faut mesurer ces grandeurs. Les clés froides restent hors ligne, sous le contrôle de plusieurs détenteurs. La compromission d'un portefeuille chaud ne fait alors perdre que cette réserve.

Vérifier les invariants en continu. La somme des soldes des utilisateurs doit être égale au total détenu sur toutes les chaînes, par actif, aux tolérances près de ce qui est en transit. Chaque ligne du grand livre doit pouvoir être remontée à un dépôt, une transaction ou un retrait. Lancer ce rapprochement selon un calendrier, et suspender automatiquement les retraits lorsqu'il échoue. Les retraits doivent pouvoir être suspendus sans interrompre le trading. Auditer aussi les soldes depuis l'extérieur, et considérer l'accès administrateur comme le principal risque.

\iqwrong{«~Conservez le solde dans le moteur d'appariement, et cela tient lieu de grand livre.~» On ne peut alors plus rapprocher les comptes, et un bogue dans l'appariement devient un bogue dans les comptes.}
\end{solution}

\begin{solution}{iq:iv3:smart-contracts-and-crypto-engineering:30}
Lisez l'état du pool et les transactions en attente depuis vos propres nœuds, au début de chaque bloc, et sur les données en attente si vous pouvez y accéder. Traitez toute transaction envoyée comme incertaine tant qu'elle n'est pas suffisamment profonde, et tenez deux chiffres distincts~: votre position confirmée et votre position attendue. Couvrez la position attendue sur les plateformes centralisées. Les coûts, les annulations et les réorganisations font tous partie du modèle.

La vue de la chaîne provient de plusieurs nœuds, car un nœud unique peut accuser du retard ou se tromper. Pour chaque bloc, il faut disposer des réserves et des ticks du pool, de vos propres soldes, des frais de base et du temps écoulé depuis le dernier bloc. Les événements de bloc doivent être indexés par le hash du bloc, afin qu'une réorganisation puisse être défaite, et il faut conserver quelques blocs d'état de retour arrière. Traitez une petite réorganisation comme un événement courant, comme vous traitez une exécution tardive sur une plateforme centralisée.

Une décision met du temps à aboutir. Une cotation calculée au début d'un créneau peut n'être incluse qu'à la fin, ou plus tard~; il faut donc inscrire dans chaque transaction une échéance et une limite de prix, de sorte qu'elle soit annulée lorsque le marché a bougé. Une annulation coûte tout de même du gas. Vous enchérissez donc pour atteindre une bonne position, et vous pouvez recourir à un canal privé pour éviter d'être pris en sandwich. Vos propres nonces exigent la même rigueur que partout ailleurs~: chaque décision conserve un seul nonce pour toujours, de sorte qu'une nouvelle tentative ne puisse pas agir deux fois, et qu'un nonce bloqué ne retienne pas tout ce qui le suit.

Une transaction envoyée passe par des états~: créée, signée, diffusée, vue dans le pool, incluse, annulée ou réussie, suffisamment profonde, et éventuellement sortie par réorganisation. La fiche de position comporte une ligne pour chacun de ces états, et se rapproche des reçus de la chaîne. Du côté de la plateforme, la même règle s'applique à ses propres exécutions. Les limites de risque s'appliquent au pire cas sur l'ensemble des éléments en attente.

Suivre l'écart entre le profit attendu et le profit réalisé pour chaque transaction, le taux d'inclusion et d'annulation, le gas consommé en part du profit, le retard des nœuds, l'inventaire sur chaque chaîne, et le glissement de la couverture. La sélection adverse se manifeste par des transactions perdantes juste après leur exécution~; il faut donc aussi la mesurer~: le flux qui échange avec un pool est souvent mieux informé.

L'environnement de test a besoin d'une chaîne locale ou d'une bifurcation~: rejouer de vrais blocs avec des réorganisations connues, annuler à la demande, et ajuster le gas. Aucun fonds réel ne doit y transiter. C'est le seul moyen sûr de vérifier qu'un retour arrière fonctionne. Aucun système de ce type n'a été exécuté ici, si bien que la conception ci-dessus n'est pas mesurée.

\iqwrong{«~Traitez la chaîne comme une plateforme de plus avec une API lente.~» Une plateforme vous indique qu'un ordre a été exécuté. Une chaîne peut revenir sur cette information pendant un certain temps.}
\end{solution}
